\documentclass[11pt,a4paper]{article}

\usepackage[T1]{fontenc}
\usepackage[utf8]{inputenc}
\usepackage{lmodern}
\usepackage[margin=2.5cm]{geometry}
\usepackage{microtype}
\usepackage{enumitem}
\usepackage{xcolor}
\usepackage{hyperref}

\setlength{\parindent}{0pt}
\setlength{\parskip}{0.6em}
\setlist[itemize]{topsep=0.2em,itemsep=0.25em}

\definecolor{linkblue}{HTML}{1F4E79}
\hypersetup{
    colorlinks=true,
    linkcolor=linkblue,
    urlcolor=linkblue,
    pdftitle={VIP Week 6 Report},
    pdfauthor={Yorgo Azzi}
}

\begin{document}

\begin{center}
    {\Large\bfseries VIP Week 6 Report}\\[0.5em]
    {\large Yorgo Azzi}\\[0.25em]
    October 21, 2026
\end{center}

\section{Overview}

This week, I integrated Google Magenta RealTime 2 with the movement-controlled
Markov music system. The objective was to use the symbolic music produced by
the existing application as live input to a generative audio model. Magenta can
then create a richer musical texture or accompaniment around either the
soprano alone or the complete four-voice output.

The complete connection was tested on a compatible machine and worked
successfully. Hand movement and Markov generation produced the symbolic input,
the input reached Magenta through its virtual MIDI port, and Magenta generated
audio in response.

\section{Work Completed}

\subsection{Magenta RealTime Integration}

Magenta RealTime 2 generates continuous stereo audio and can be controlled by
a text prompt and live MIDI notes. The text prompt determines the general
instrumentation and musical character, while the MIDI input provides pitch,
timing, note onset, and sustain information.

The output of the existing system is captured after the Markov engine and the
hand-controlled soprano have been combined. Each generated event contains the
current pitches, new note onsets, released notes, sustained notes, velocity,
and timing. These events are converted into the 128-pitch MIDI-conditioning
state used by Magenta, where each pitch is represented as off, sustained, or a
new onset.

The conditioning data are updated at 25 frames per second. This preserves short
notes and held notes without making Magenta wait for the next measure or chord
change. Pausing playback, disabling Magenta, or switching input modes sends a
clear note reset so that notes do not remain stuck.

\subsection{Soprano and Full-Texture Inputs}

I added two selectable Magenta input modes:

\begin{itemize}
    \item \textbf{Soprano only} sends only the final soprano line. In Full SATB
    mode, this soprano comes from the Markov and voice-leading engine. In
    Soprano Performance mode, it comes from the performer's right-hand
    movement. Magenta is then free to infer an accompaniment around the melody.
    \item \textbf{All four voices} sends soprano, alto, tenor, and bass. This
    gives Magenta the complete harmony and rhythm produced by the Markov system,
    so its generated audio can follow the existing musical structure more
    closely.
\end{itemize}

The camera application now includes a Magenta On/Off button, an input selector,
and a connection-status display. These controls can be changed during a
performance without restarting the hand tracker or Markov generator.

\subsection{Local Gateway and Testing Support}

The first implementation uses a local gateway between the camera application
and Magenta. The realtime application converts the generated music into a
fixed-rate conditioning stream and sends it over the computer's loopback
address. The gateway then forwards the note states as MIDI messages to the
virtual \texttt{Magenta RT Input} port created by the official Magenta
standalone application.

This design also includes a mock backend. The mock receives and displays the
same pitches without running the Magenta neural model, allowing the complete
data path, filtering, and interface status to be tested on the development
laptop. Communication failure does not stop the camera, Markov model, MIDI
playback, or live score.

\section{Results}

The main results from this week were:

\begin{itemize}
    \item The existing hand and Markov performance can now be used as live MIDI
    conditioning for Magenta RealTime 2.
    \item The user can switch between soprano-only and complete SATB input from
    the integrated application.
    \item The final performed soprano is used, so Magenta receives the Markov
    soprano in Full SATB mode and the hand-controlled soprano in Soprano
    Performance mode.
    \item Onsets, sustained notes, releases, repeated notes, short notes, pause,
    and mode changes are represented correctly in the conditioning stream.
    \item A mock backend allows the connection to be verified without loading
    the Magenta model.
    \item The full system was tested on a compatible machine. The gateway found
    the Magenta virtual MIDI input, the interface reported a successful
    connection, and Magenta produced audio from the generated musical input.
\end{itemize}

\section{Challenges}

Magenta produces stereo audio rather than symbolic MIDI output. Its generated
notes therefore do not appear in the live score, which continues to show only
the music created by the Markov and hand-performance system. The output is also
influenced by Magenta's text prompt and internal generation, so it is not a
deterministic orchestration of the input notes. Soprano-only input gives the
model more freedom, while all-voice input constrains it more strongly and may
produce a denser result.

The Magenta audio also has more latency than the direct MIDI soprano because it
must pass through model inference and audio buffering. Keeping the immediate
soprano sound alongside the generated background may therefore provide a
stronger feeling of direct control than using only Magenta audio.

Finally, the current local protocol and gateway work, but they add an extra
process and communication layer. Since the official Magenta application
already provides a virtual MIDI input on the same computer, the performance
application could send MIDI to it directly. This would preserve the musical
behavior while simplifying setup and reducing the number of files and commands.

\section{Next Steps}

\begin{itemize}
    \item Simplify the architecture by sending the selected soprano or SATB
    notes directly to \texttt{Magenta RT Input}. Keep the mock option inside the
    performance application and remove the additional local protocol and
    gateway process.
    \item Compare soprano-only and all-voice conditioning using the same prompts
    and musical gestures. Evaluate harmonic agreement, responsiveness, musical
    density, and the amount of creative freedom given to Magenta.
    \item Replace the traditional sheet-music view of the soprano with a live
    piano-roll display. This should make the performed pitch direction, note
    duration, and timing easier to follow during fast hand-controlled passages.
\end{itemize}

\section{Resources}

\begin{itemize}
    \item Project repository:
    \url{https://github.com/YorgoAzzi/MarkovModel}
    \item Magenta RealTime 2 repository:
    \url{https://github.com/magenta/magenta-realtime}
    \item Magenta RealTime 2 model description:
    \url{https://github.com/magenta/magenta-realtime/blob/main/MODEL.md}
    \item Magenta standalone application documentation:
    \url{https://github.com/magenta/magenta-realtime/blob/main/docs/apps/standalone_app.md}
\end{itemize}

\end{document}
